\chapter{Literature Review}

The literature review is structured into three main themes relevant to this project: the ecological necessity of coral reef monitoring, the evolution of underwater scene change detection methods, and the recent rapid development of underwater 3D Gaussian splatting techniques.

\section{The Importance of Coral Monitoring and Traditional Methods}

Coral reef monitoring forms the fundamental baseline for marine conservation, providing the critical data necessary to detect long term ecological trends, evaluate management interventions, and provide early warning systems for disturbances such as mass bleaching and disease outbreaks~\cite{souter2020status}. Traditional monitoring methodologies predominantly rely on diver-based in-situ surveys, utilizing techniques such as Manta Tows for broad-scale assessment, Point Intercept Transects (PIT) for benthic composition, and fixed-site photo-quadrats for localized structural tracking~\cite{icri2021monitoring}. However, as highlighted by a recent comprehensive review of coral reef condition indicators by Davey et al.~\cite{davey2025review}, there is a common disconnect between the vast amounts of monitoring data collected globally and its actual application in spatial conservation planning. Traditional 2D metrics like percentage live coral cover, while ubiquitous, often fail to capture the complex, 3D structural dynamics necessary for holistic ecosystem assessment~\cite{davey2025review}. This limitation underscores the urgent need for advanced, scalable, and structurally aware monitoring technologies such as 3D Gaussian Splatting that can bridge the gap between data collection and actionable conservation insights.

\section{Scene Change Detection and 3D Representations}

\textbf{Render-then-Compare Paradigms vs. Primitive-Level Comparison.} 
Traditional scene change detection (SCD) across unconstrained viewpoints has historically operated in image space, relying on learned features or foundation models. Recent methods leverage 3D Gaussian Splatting (3DGS)~\cite{10897713} to synthesize reference views at novel poses, scoring changes by rendering views and comparing features or pixels~\cite{10.1145/3806262.3806264}. While these render-then-compare approaches successfully suppress  view inconsistent distractors, they fundamentally use 3D primitives merely as a rendering substrate, relying on 2D space for the actual change signals. This often struggles to explicitly separate structural changes from mere illumination or appearance shifts.

In contrast, recent advancements have introduced a new paradigm operating directly within the 3D representation space. GS-DIFF~\cite{galappaththige2026pixelsprimitivesscenechange} pioneered direct primitive-space comparison by analyzing the core attributes of 3D Gaussians (position, covariance, and color). The primary strength of this approach is its ability to separate structural transformations from superficial appearance shifts without relying on 2D image synthesis or external supervision. Similarly, 3DGS-CD~\cite{10852207} leverages 2D segmentation priors fused into the 3D space to detect physical object rearrangements rapidly, without requiring depth inputs or pre-defined object models. However, the critical limitation of primitive-level comparison is the inherently under-constrained nature of 3DGS; independent optimization runs can yield drastically different primitive configurations (Gaussian drift) even for unchanged scenes. This necessitates complex modeling of geometric and photometric drift to prevent false positives~\cite{galappaththige2026pixelsprimitivesscenechange}.

\textbf{Spatial Consistency Constraints.}
To address the spatial inconsistency often found in localized change masks, methods like SCAR-3D~\cite{zhou20253dscenechangemodeling} propose multi-view aggregation frameworks. By employing signed-distance-based differencing followed by voting and pruning mechanisms, this approach robustly separates pre- and post-change states. While highly accurate in producing consistent change masks, exploring non-rigid deformations within this framework remains an ongoing challenge.

\section{Underwater Scene Reconstruction and 3D Gaussian Splatting}

Underwater image formation is heavily degraded by absorption, scattering, and color shift, historically addressed via single-image enhancement~\cite{s20185170} or Neural Radiance Fields~\cite{10203316}. With the advent of 3DGS, the literature has diverged into three primary methodological mindsets to tackle these aquatic challenges.

\textbf{Approach 1: Physically-Grounded Image Formation (Inverse Rendering).} 
Several works directly embed mathematical models of light attenuation and scattering into the splatting optimization pipeline. SeaSplat~\cite{11128502} constrains 3DGS using range and color dependent models, while WaterGS~\cite{article} formulates reconstruction as a formal inverse rendering problem. UW-3DGS~\cite{xing2025uw3dgsunderwater3dreconstruction} employs a learnable image-formation module paired with Physics-Aware Uncertainty Pruning, and DualPhys-GS~\cite{LI2025104405} introduces dual feature-guided scattering-attenuation models. RecGS~\cite{zhang2024recgsremovingwatercaustic} specifically targets physical artifacts like water caustics. 
\\ \textit{Strengths and Limitations:} The core advantage of this physical mindset is its ability to disentangle true object color and geometry from the water medium, yielding highly accurate, structurally consistent primitives. However, this comes at a significant computational cost. Furthermore, these methods are highly sensitive to the accuracy of depth priors (which are notoriously difficult to obtain underwater) and often struggle in extreme long-distance scenes where the scattering signal overwhelmingly dominates the object signal.

\textbf{Approach 2: Prior-Driven and Image Restoration Fusion.}
An alternative mindset avoids explicit physical equations in favor of leveraging pre-enhanced views or data-driven image restoration priors. Fast methods like WaterSplatting~\cite{li2025watersplattingfastunderwater3d} and R-Splatting~\cite{huang2025restorationreconstructionrethinking3d} bridge traditional 2D underwater restoration networks with the 3D structural constraints of Gaussian splatting. WaterClear-GS~\cite{zhang2026watercleargsopticalawaregaussiansplatting} uses optical-aware dual-branch optimization to clean textures. 
\\ \textit{Strengths and Limitations:} By optimizing for visual appearance rather than strict physics, these methods train significantly faster and frequently achieve superior perceptual quality metrics (e.g., PSNR, SSIM). Unfortunately they are highly susceptible to view inconsistencies introduced by 2D filters, are prone to overfitting localized noise, and inherently lack the capability to model distant volumetric media correctly since they rely entirely on the explicit geometric representations of 3DGS.

\textbf{Approach 3: Multi-Modal Sonar Fusion.}
To overcome the severe visual degradation in turbid waters, a growing body of work integrates acoustic sensor data. SonarSplat~\cite{11223217} and NAS-GS~\cite{xu2026nasgsnoiseawaresonargaussian} synthesize novel sonar views while modeling acoustic streaking phenomena. Crucially, Aqua-Splat~\cite{11184160} and SonarReg-GS~\cite{s26154713} fuse Forward Looking Sonar (FLS) measurements to serve as sparse metric depth anchors.
\\ \textit{Strengths and Limitations:} Sonar effectively penetrates highly turbid water, and fusing it with 3DGS elegantly solves the monocular depth scale ambiguity problem that plagues purely visual underwater SLAM. Conversely, the limitations are severe: raw sonar observations are highly sparse, noisy, and suffer from elevation angle ambiguities. While sonar fusion corrects global scale, it struggles to provide the dense geometric recovery required for fine-grained structural modeling without complex multi-sensor calibration.

\textbf{Geometric Adaptations for Extreme Environments.} 
Finally, to survive dynamic elements and low-visibility, some methods directly alter the structural mechanics of the Gaussians. Water-Adapted 3DGS~\cite{10.3389/fmars.2025.1573612} introduces complexity-adaptive point distributions, while RUSplatting~\cite{jiang2025rusplattingrobust3dgaussian} uses frame interpolation for sparse views. Dynamic distractors (e.g., marine life) are filtered by UW-GS~\cite{10943785} and Spatiotemporal 3DGS~\cite{Liu_2025}. For wide-field observation, Underwater360~\cite{hu2026underwater360reconstructingunderwaterscenes} leverages omnidirectional splatting. Controlled cross regime studies~\cite{alvarez2026gaussiansplattingunderwatercontrolled} and surface-aligned mesh extractions like SuGaR~\cite{guedon2024sugar} further validate the fidelity of these environmental adaptations.

\section{Summary of Findings and Research Gap}
The literature demonstrates a clear trajectory, the field has rapidly moved from adapting basic 3DGS to underwater scenes via standard image restoration priors, to embedding complex physical light-attenuation models directly into the optimization, and finally to adapting the structural primitives themselves to survive extreme turbidity and integrate multi-modal sensors. For coral reef ecosystems which present highly complex geometries, transient marine life, and dynamic lighting conditions these specialized underwater adaptations are paramount. 

Table~\ref{tab:3dgs_methods} provides a consolidated overview of the various 3DGS methods and their primary contributions to scene change detection and underwater reconstruction.

\begin{table}[htbp]
    \centering
    \caption{Summary of 3DGS Methods and Approaches}
    \label{tab:3dgs_methods}
    \resizebox{\textwidth}{!}{
    \begin{tabular}{@{} l l p{8.5cm} @{}}
    \toprule
    \textbf{Category} & \textbf{Method} & \textbf{Key Approach / Contribution} \\ \midrule
    \multirow{2}{*}{\parbox{3.5cm}{\textbf{Scene Change\\ Detection}}} 
    & GS-DIFF~\cite{galappaththige2026pixelsprimitivesscenechange} & Direct primitive-space comparison (position, covariance, color). \\
    & 3DGS-CD~\cite{10852207} & Detects physical object rearrangements using explicit splat geometric properties. \\ \midrule
    \multirow{5}{*}{\parbox{3.5cm}{\textbf{Physically-Grounded\\ Underwater 3DGS}}} 
    & SeaSplat~\cite{11128502} & Physically grounded range and color dependent modeling. \\
    & WaterGS~\cite{article} & Formulates scene reconstruction as an inverse rendering problem. \\
    & DualPhys-GS~\cite{LI2025104405} & Dual feature-guided scattering-attenuation modeling for selective wavelength attenuation. \\
    & UW-3DGS~\cite{xing2025uw3dgsunderwater3dreconstruction} & Learnable image-formation with Physics-Aware Uncertainty Pruning. \\
    & RecGS~\cite{zhang2024recgsremovingwatercaustic} & Targets and removes specific physical artifacts like water caustics. \\ \midrule
    \multirow{2}{*}{\parbox{3.5cm}{\textbf{Image Restoration\\ \textit{\&} Prior-Driven}}} 
    & WaterClear-GS~\cite{zhang2026watercleargsopticalawaregaussiansplatting} & Optical-aware dual-branch optimization using restoration priors. \\
    & Fast Methods~\cite{li2025watersplattingfastunderwater3d, huang2025restorationreconstructionrethinking3d} & Bridges traditional 2D underwater restoration networks with 3D structural constraints. \\ \midrule
    \multirow{5}{*}{\parbox{3.5cm}{\textbf{Extreme Environments\\ \textit{\&} Adaptations}}} 
    & Water-Adapted 3DGS~\cite{10.3389/fmars.2025.1573612} & Complexity-adaptive points and depth-adaptive multi-scale rendering. \\
    & RUSplatting~\cite{jiang2025rusplattingrobust3dgaussian} & Frame interpolation and decoupled RGB learning for sparse-view/low-light. \\
    & UW-GS~\cite{10943785} & Distractor-aware variance filtering for transient marine life. \\
    & Spatiotemporal 3DGS~\cite{Liu_2025} & Models localized turbidity and current variations. \\
    & Underwater360~\cite{hu2026underwater360reconstructingunderwaterscenes} & Omnidirectional splatting to eliminate refraction distortion. \\ \midrule
    \multirow{2}{*}{\parbox{3.5cm}{\textbf{Multi-Modal\\ Sonar Fusion}}} 
    & SonarSplat~\cite{11223217}, NAS-GS~\cite{xu2026nasgsnoiseawaresonargaussian} & Synthesizes novel sonar views and handles noise-aware splatting.\\
    & Aqua-Splat~\cite{11184160}, SonarReg-GS~\cite{s26154713} & Fuses sonar for sparse metric depth anchors to correct scale ambiguity. \\ \bottomrule
    \end{tabular}
    }
\end{table}

This project aims to bridge the critical gap between advanced underwater 3DGS reconstruction (Section 2.3) and primitive-space scene change detection (Section 2.2). By evaluating which of the state-of-the-art underwater 3DGS methods produces the most geometrically robust and physically accurate primitives for coral scenes, this project will establish a foundational pipeline for tracking structural reef changes. Ultimately, this will move coral monitoring beyond simple 2D metrics (Section 2.1) and provide the rich, 3D structural data necessary to support effective spatial conservation planning.
